\subsection{平均加速、总周期与执行效率的不同视角}
表\ref{tab:q1-main}中的平均加速比按式\eqref{eq:mean-speedup-derivation}逐图等权统计。为了评价这批图的总体执行负担，表\ref{tab:q1-aggregate-perspective}另把单核总周期 $312\,479\,373$ 作为分母，列出合计周期缩减率与总周期比。数值均直接由同一500行结果表计算，不引入另一批算例。五核总周期较单核下降 $76.5839\%$，其倒数形式的总周期比为 $4.270565$；规定的逐图平均加速比却为 $3.665903$。差异并非评估矛盾，而是式\eqref{eq:weighted-speedup}的权重差异。论文主图仍采用规定的逐图等权均值，总周期比只用于说明大图的绝对时延贡献。
\begin{table}[!htbp]\centering
\caption{问题一结构感知搜索百图总周期视角与逐图等权视角}\label{tab:q1-aggregate-perspective}
\begin{tabular}{rrrrr}
\toprule
核数 & 总周期比 & 总周期下降率 & 平均加速比 & 平均并行效率\\
\midrule
1 & 1.000000 & $0.0000\%$ & 1.000000 & 1.000000\\
2 & 1.968419 & $49.1978\%$ & 1.834917 & 0.917459\\
3 & 2.843802 & $64.8358\%$ & 2.546238 & 0.848746\\
4 & 3.619570 & $72.3724\%$ & 3.161009 & 0.790252\\
5 & 4.270565 & $76.5839\%$ & 3.665903 & 0.733181\\
\bottomrule
\end{tabular}
\end{table}
并行效率定义为 $\eta_k=\overline S_k/k$，只是把相对核数的等权加速程度标准化。五核时 $\eta_5=0.733181$ 并不意味着其余 $26.6819\%$ 的周期都由 DDR 独自造成；关键链、Pipe 不均、释放等待和核内溢出都可能参与。任务执行时间越短是首要目标，不能为了让 $\eta_k$ 的比值好看而故意牺牲 Makespan。后文图\ref{fig:q1-total}的两个子图分别展示逐图平均加速比与百图总周期，报告两个互补视角。


\subsection{典型结果与反例}
第59图的单核参考为 $829\,473$ 周期，五核方案为 $150\,782$ 周期，逐图加速比约 $5.5011$；其额外搬运则从单核 $1\,691\,648$ 增至五核 $4\,325\,376$ 字节。这个例子说明完成时间改善可以伴随更多逻辑搬运，也说明超过五倍的比值并不等于突破硬件计算吞吐上限。赛题的单核分母来自固定的核内启发式，不是单核全局最优；多核切图改变核内问题规模和溢出行为，可能使官方启发式在五核下比把原单核结果简单除以五更有效。第97图单核 $11\,491\,509$、三核 $3\,509\,566$ 周期，加速比约 $3.2743$，是三核配置的分布端点之一。端点代表有利结构，不代表百图平均水平，也不能据此承诺每张图接近理想线性加速。

第6图呈现相反的警示：三核 $32\,925$ 周期、四核 $31\,129$ 周期、五核 $31\,729$ 周期，五核比四核慢。其核数、划分与映射共同改变，不能仅凭这三个总数断言“多一核导致 DDR 拥塞”；需结合时间线中的关键依赖与在途传输核查。该图在前后算法版本比较中三核和四核亦是两处退步，说明有限预算下的新候选机制可能挤掉旧版有价值的候选。故当前选优只保证\emph{同一候选池内}不接受实测变差方案，跨版本必须再做同图同核配置对照。

\subsection{额外搬运与大图求解代价}
二核百图总额外搬运量为 $1\,972\,107\,414$ 字节，高于三核、四核和五核的合计值；五核也略高于四核。这一序列不单调，表明“核数增加必然增加割边搬运”过于简单。算例的实际划分粒度、完整分量分配、共享输入重复读取和核内溢出都随核数变动；如果把所有差值都归咎于跨核边数，便忽略了场景 A 下\emph{同核跨 Task 仍须 DDR}的规则。对本文采用方案，搬运量是次级实测指标，不能由算子数或切边数推算为附录最终值。

大图还有另一种与 NPU 周期不同的成本：主机求解墙钟时间。官方核内启发式、合法性检查和全局事件评估对大型图更贵，若每生成一个局部候选就重新准备全部核内任务，求解过程会浪费大部分时间。对已有结构使用哈希键复用只读准备结果、对商图改变区域增量检查无环，减少的是\emph{找出方案的主机计算}；它不改变候选的官方 $T^A(p)$。本轮100图多核热搜索记录的最长时长约 $363.859$ 秒，含义仅限该阶段；冷启动单核基准和独立最终回放需另外计时。因此本文没有把“每张图端到端5--10分钟”作为已通过的全量实测结论。
